\documentclass[10pt,a4paper]{amsart}
\usepackage[margin=19mm]{geometry}
\usepackage{amsmath,amssymb,mathtools}
\usepackage[scr=rsfs,cal=euler]{mathalfa}
\usepackage{xfrac}
\usepackage[unicode,hidelinks,bookmarks=false]{hyperref}
\newcommand{\ie}{i.e.\ }
\newcommand{\cf}{cf.\ }
\newcommand{\resp}{respectively\ }
\newcommand{\Subsection}{\subsection}
\providecommand{\nobreakdash}{\nobreak\hbox{-}}
\setlength{\emergencystretch}{2em}
\multlinegap0pt
\usepackage{amssymb}
\newtheorem{lemma}{Lemma}
\newtheorem{proposition}{Proposition}
\newcommand {\Ga}{\Gamma}
\newcommand {\bCP} {\mathbb {CP}}
\title{Ribbon graphs and meromorphic functions}
\author{Boris Shapiro}
\date{Source: arXiv:2604.21358v1 (CC BY 4.0)}
\begin{document}
\maketitle

\noindent\emph{Source and licence.} Ribbon graphs and meromorphic functions. Version arXiv:2604.21358v1 (CC BY 4.0), distributed by arXiv under the Creative Commons Attribution 4.0 International licence, http://creativecommons.org/licenses/by/4.0/, as recorded in the arXiv licence metadata for this exact version and shown on the abstract page https://arxiv.org/abs/2604.21358v1. Full licence notice: \texttt{licenses/arxiv-2604.21358-CC-BY-4.0.txt}.\par\smallskip

\noindent\textit{Editorial scope.} Counted unit: the article's proposition pr:im (source0.tex lines 241--254). The article's complete original proof of it is delivered with it (source0.tex lines 256--272, one \texttt{proof} environment of 17 lines). No mathematics is written, repaired or re-derived: the statement and the proof are the article's own lines, in the article's own notation, and no retained line is substituted or re-flowed.  Retained uncounted as the source-local context the proof needs: lines 223--225.  Nothing was omitted from any retained span, so the package carries no omission record.  No retained line contains a citation, so the wrapper carries no bibliography and no bibliography entry.  No retained line contains a figure environment, an image-inclusion command or drawing code, so no image is delivered and the package declares no asset dependency.  Editorial formatting only, in addition: an AMS \texttt{amsart} preamble that restates the article's own declarations and macros used by the retained lines, namely the environments \texttt{lemma}, \texttt{proposition} on separate counters, together with the standard packages those lines need.  The retained environments print in this document as Lemma 1, Proposition 1 in order of appearance, whereas the article prints the same items with the numbering of its own full document.  The complete native source of the article as distributed in the arXiv e-print for this exact version is bundled unchanged as \texttt{sources/199020/source0.tex}.
\begin{lemma}\label{lm:noselfint}
For any class $[\Ga]$, there exists a representative $\widetilde \Ga\in[\Ga]$ such that $\phi(\widetilde \Ga)$ has no self-intersections of individual edges. Furthermore, any representative $\Ga^*\in[\Ga]$ realizing the minimum $\sharp([\Ga],\phi)$ has this property.
\end{lemma}

\begin{proposition}\label{pr:im}
\rm
(i) For any compact Riemann surface $Y$ of genus $g$, any class of ribbon graphs $[\Ga]$, and any meromorphic function $\phi:Y\to \bCP^1$,
\[
\sharp([\Ga],\phi)\ge g.
\]

\smallskip
\noindent
(ii) For any rational function $\phi:\bCP^1\to \bCP^1$ and any class $[\Ga]$,
\[
\sharp([\Ga],\phi)=0.
\]
\end{proposition}

\begin{proof}[Proof of Proposition~\ref{pr:im}]
(i) Choose a representative $\Ga^*\in[\Ga]$ realizing $\sharp([\Ga],\phi)$ and, by Lemma~\ref{lm:noselfint}, assume that no individual edge has self-intersections in the image. Let
\[
d=\sharp([\Ga],\phi)=\sharp(\Ga^*,\phi).
\]
Resolve each transverse double point of the immersed graph $\phi(\Ga^*)\subset \bCP^1$ by inserting a $4$-valent vertex. This produces an embedded graph $H\subset \bCP^1$. Let $N(H)$ be a closed ribbon neighbourhood of $H$ in the sphere. Since $H$ is embedded in $\bCP^1$, the surface obtained from $N(H)$ by capping its boundary components with disks has genus $0$.

Now reconstruct the original immersed image from $H$ one crossing at a time. At a former double point, the passage from the neighbourhood of the $4$-valent vertex in $N(H)$ to the neighbourhood of a transverse crossing is achieved by attaching one $1$-handle to the thickened surface. Consequently, after all $d$ crossings are restored, the closed surface determined by the resulting ribbon graph has genus at most $d$.

On the other hand, the ribbon surface determined by $\Ga^*$ is precisely $Y$, because $\Ga^*$ spans $Y$. Hence
\[
g(Y)=g\le d=\sharp([\Ga],\phi),
\]
which proves (i).

(ii) Let $Y=\bCP^1$ and let $\phi:\bCP^1\to \bCP^1$ be rational. Choose a closed topological disk $D\subset \bCP^1$ avoiding the critical points of $\phi$ and small enough that $\phi|_D$ is injective. Since every orientation-preserving homeomorphism of the sphere is homotopic to the identity, any embedded graph on $\bCP^1$ can be moved by such a homeomorphism into $D$. Hence every class $[\Ga]$ contains a representative $\Ga'\subset D$. For this representative, $\phi(\Ga')$ is embedded, so $\sharp(\Ga',\phi)=0$. Therefore $\sharp([\Ga],\phi)=0$.
\end{proof}

\end{document}
